SEGLEA computes the source effective-address point and applies the selected segment's pre-translation to that point. The size suffix selects the index scale used during effective-address calculation. On success, SEGLEA writes the segment result to the destination and clears FLAGS. A failed segment-bound check sets V while clearing Z, N, and C, suppresses the destination write, commits effective-address auto-update effects, and completes without an exception. REPcc observes the failure result as a B-sized 0 or 1 rather than architectural FLAGS.

Let \(a\) be the source effective address before a memory read. The
\hyperref[page:effective-addressing-modes]{effective-address form} selects the segment image, which supplies the
\(\mathit{base}\), \(\mathit{span}\), and \(\mathit{limit}\) quantities defined in
\hyperref[section:segment-registers]{Segment Registers}.
A zero mantissa disables bounds and returns \(a\). In translated mode, \(a\) must satisfy
\(0 \le a < \mathit{span}\) and the result is \(\mathit{base}+a\). In bounds-only mode, \(a\) must satisfy
\(\mathit{base} \le a < \mathit{limit}\) and the result remains \(a\). An out-of-bounds address writes FLAGS as
Z=N=C=0 and V=1, in which case the destination is not written. A successful address writes FLAGS as Z=N=C=V=0.
